\begin{abstract}
We introduce \textsc{QLoRA}, a resource-efficient finetuning strategy that substantially lowers memory consumption, enabling the finetuning of models with up to 65B parameters on a single 48GB GPU without sacrificing the effectiveness of standard 16-bit finetuning tasks. In \textsc{QLoRA}, gradients are propagated through a static, 4-bit quantized pretrained language model into Low Rank Adapters (LoRA). The top-performing suite of models from our work, dubbed \textbf{Guanaco}, achieves superior results compared to all previously available open models on the Vicuna benchmark, attaining 99.3\% of ChatGPT's performance, and requires just 24 hours of training on one GPU. \textsc{QLoRA} integrates several technical advancements to optimize memory utilization without diminishing performance: (a) 4-bit NormalFloat (NF4), a novel data format tailored for normal distributions of weights to achieve information-theoretical optimality; (b) Double Quantization, which further reduces memory requirements by applying quantization to the quantization coefficients; and (c) Paged Optimizers, which help regulate transient memory increases. Employing \textsc{QLoRA}, we finetuned over 1,000 models, conducting a comprehensive examination of instruction-following and conversational abilities across 8 instruction datasets, various architectures (LLaMA, T5), and larger model sizes—including 33B and 65B parameters—that are impractical with conventional finetuning approaches. Our findings demonstrate that \textsc{QLoRA} finetuning with concise, high-quality datasets achieves state-of-the-art performance, even for models smaller than the current state-of-the-art. We present an in-depth assessment of chatbot behavior using both human and GPT-4 based ratings, establishing that GPT-4 evaluation offers a practical and low-cost alternative to human judgment. Additionally, we observe that existing chatbot benchmarks are unreliable indicators of chatbot proficiency. Through a targeted "lemon-picked" analysis, we illustrate cases where \textbf{Guanaco} underperforms relative to ChatGPT. All models and source code, including CUDA kernels for 4-bit training, are openly available.\footnote{\url{https://github.com/artidoro/qlora} and \url{https://github.com/TimDettmers/bitsandbytes}}
\end{abstract}

\section{Introduction}

The process of finetuning large language models (LLMs) stands as a robust approach for enhancing their capabilities \citep{min2021metaicl, wei2021finetuned, ouyang2022training, wang2022super, wang2022self, liu2022few} and for shaping model behaviors by introducing or eliminating specific characteristics \citep{ouyang2022training,askell2021general,bai2022training}. Yet, adapting extremely large models via finetuning presents significant financial and computational challenges; for instance, standard 16-bit finetuning of the 65B-parameter LLaMA model~\citep{touvron2023llama} consumes upwards of 780 GB of GPU memory. Though advances in quantization have managed to compress LLMs for inference \citep{dettmers2022llm, dettmers2022case, frantar2022gptq, xiao2022smoothquant}, these solutions fail to extend to the training phase, where they become unstable \citep{wortsman2023stable}.

In this work, we provide the first evidence that it is feasible to finetune a 4-bit quantized model without experiencing any loss in accuracy. Our proposed approach, \textsc{QLoRA}, leverages an innovative high-precision quantization method that converts a pretrained model to 4 bits and incorporates a compact set of trainable Low-rank Adapter weights \citep{hu2021lora}

which are adjusted via backpropagation through the quantized parameter space.

With \textsc{QLoRA}, memory demands for finetuning a 65B parameter model drop from over 780GB to less than 48GB, maintaining both speed and accuracy relative to the 16-bit finetuned baseline. This advancement greatly broadens the accessibility of LLM finetuning—enabling single-GPU adaptation for even the largest open models to date. Leveraging \textsc{QLoRA}, we introduce the \textbf{Guanaco} model family; notably, our second-best model attains 97.8\% of ChatGPT's performance on the Vicuna~\citep{vicuna2023} evaluation while completing training in under 12 hours on a single consumer GPU. By utilizing a single professional GPU for 24 hours, our largest Guanaco model achieves 99.3\%, effectively matching ChatGPT’s performance on Vicuna. Furthermore, upon deployment, our smallest \textbf{Guanaco} variant (7B parameters) operates with only 5 GB of memory and surpasses the 26 GB Alpaca model by more than 20 points on the Vicuna benchmark (see Table~\ref{tbl:vicuna_eval}).

\textsc{QLoRA} incorporates several key innovations that significantly decrease memory consumption while maintaining model performance: (1) \textbf{4-bit NormalFloat}, a quantization datatype that is theoretically optimal for data with a normal distribution, demonstrating empirically superior performance to both 4-bit Integers and 4-bit Floats. (2) \textbf{Double Quantization}, a technique that further compresses by quantizing the quantization constants themselves, resulting in an average memory saving of around 0.37 bits per parameter (equivalent to roughly 3 GB for a 65B-parameter model). (3) \textbf{Paged Optimizers}, which leverage NVIDIA unified memory to circumvent memory spikes caused by gradient checkpointing when processing mini-batches with long sequence lengths. These components are integrated within an enhanced LoRA methodology, which applies adapters at each network layer and thus nearly eliminates the accuracy compromises that prior approaches encountered.

The high efficiency of \textsc{QLoRA} makes it feasible to systematically investigate instruction finetuning and chatbot capabilities across a range of model sizes that would be otherwise prohibitive due to memory constraints associated with traditional finetuning. To this end, we trained over 1,000 distinct models spanning multiple instruction tuning corpora, model architectures, and parameter counts from 80M to 65B. Beyond demonstrating that \textsc{QLoRA} achieves parity with 16-bit training performance (\S\ref{sec:comp_to_std}) and is capable of powering a state-of-the-art chatbot, \textbf{Guanaco} (\S\ref{sec:pushing}), we examine patterns among trained models. Our analyses reveal, first, that the quality of training data outweighs the importance of data quantity—for instance, the 9k-sample OASST1 dataset led to superior chatbot outcomes compared to the much larger 450k-sample (subsampled) FLAN v2 dataset, despite both being curated for instruction-following tasks. Secondly, our results indicate that excelling on the Massive Multitask Language Understanding (MMLU) benchmark does not guarantee competitive performance on the Vicuna chatbot benchmark, and vice versa; in essence, selecting the appropriate dataset for the target task has a greater impact than simply scaling dataset size.

Additionally, we conduct a thorough examination of chatbot effectiveness utilizing assessments from both human evaluators and GPT-4. Our evaluation employs a tournament-style framework where models face off to generate optimal responses to selected prompts. The victor of each contest is determined by either GPT-4 or by human judges. We summarize the outcomes of these tournaments by calculating Elo scores~\citep{elo1967proposed, elo1978rating}, which serve as the basis for establishing the performance hierarchy of the chatbots. Our findings reveal that rankings generated by GPT-4 generally concur with those from human evaluators, though there are notable instances of disagreement. Consequently, we underscore that while automated, model-based evaluation offers a cost-efficient alternative to manual annotation, it is accompanied by certain uncertainties.

To complement our quantitative chatbot benchmarks, we perform an in-depth qualitative assessment of the \textbf{Guanaco} models. This assessment uncovers both exemplary and problematic outputs that may not be reflected in purely quantitative metrics.

For transparency and further research, we provide all model-generated outputs together with both human and GPT-4 assessments. Our entire codebase and CUDA kernels are released as open-source resources, and our techniques have been incorporated into the Hugging Face transformers repository \citep{wolf2019huggingface} for broad accessibility. Furthermore, we offer a suite of adapters covering models of sizes 7/13/33/65B, each trained on eight distinct instruction-following datasets, yielding a collection of 32 openly available, fine-tuned models.

\section{Background}
\paragraph{Block-wise k-bit Quantization} Quantization refers to transforming input data from a format with greater informational granularity to one with lower granularity, typically by reducing the number of bits used to represent each value (for example, converting from 32-bit floating point numbers to 8-bit integers). To effectively utilize the dynamic range of the reduced-bit format, data are typically rescaled so that the range aligns with that of the target data type. This normalization typically employs the absolute maximum value present in the input tensor. For instance, to quantize a tensor of 32-bit floating point values (FP32) to an 8-bit integer (Int8) within the interval $[-127, 127]$:

\begin{equation}
    \mathbf{X}^{ \text{Int8}} = \text{round}\left( \frac{127}{{\text{absmax}}(\mathbf{X}^{{\text{FP32}}})}\mathbf{X}^{{\text{FP32}}} \right) = \text{round}(c^{\text{FP32}}\cdot \mathbf{X}^{{\text{FP32}}}),
\end{equation}

with $c$ representing the {\em quantization constant} or {\em quantization scale}. The dequantization process, which reverses quantization, is expressed as:

\begin{equation}
    \text{dequant}(c^{\text{FP32}}, \mathbf{X}^{\text{Int8}}) = \frac{\mathbf{X}^{\text{Int8}}}{c^{\text{FP32}}} = \mathbf{X}^{\text{FP32}}
\end{equation}

A significant drawback of this method arises when the input tensor contains a value of unusually large magnitude (an outlier). In such cases, many of the quantization bins—particular combinations of bits—are poorly populated, with few or even no elements being assigned to certain bins, resulting in inefficient usage. A widely-adopted solution to address the issue of outliers is to partition the input tensor into separate blocks, each of which is quantized independently and assigned its own quantization constant $c$. Specifically, the input tensor $\mathbf{X}\in \mathbb{R}^{b\times h}$ is reshaped into $n$ sequential blocks of size $B$ by flattening it and dividing the resulting sequence into $n = ({b\times h} )/{B}$ segments. Each block is then quantized in isolation, following Equation~1, thereby yielding a quantized tensor and an associated set of $n$ quantization constants $c_i$.

\paragraph{Low-rank Adapters} LoRA (Low-rank Adapter) finetuning \citep{hu2021lora} is a technique aimed at minimizing memory consumption by employing a small collection of trainable parameters—referred to as adapters—while the original model weights remain frozen throughout training. During the optimization process, gradients computed by stochastic gradient descent traverse the fixed pretrained parameters, updating only the adapter components to minimize the chosen loss function. LoRA enhances a linear projection layer by introducing an additional factorized (low-rank) projection. For a linear operation defined as ${\mathbf{X} \mathbf{W} = \mathbf{Y}}$ with $\mathbf{X}\in  \mathbb{R}^{b\times h}$ and $\mathbf{W}\in  \mathbb{R}^{h\times o}$, LoRA produces:

\begin{equation}
    \mathbf{Y} = \mathbf{X}\mathbf{W} + s\mathbf{X}\mathbf{L}_1\mathbf{L}_2,
\end{equation}

in which $\mathbf{L}_1\in  \mathbb{R}^{h\times r}$ and $\mathbf{L}_2\in  \mathbb{R}^{r\times o}$ are trainable low-rank matrices, and $s$ denotes a scalar value.

\paragraph{Memory Requirement of Parameter-Efficient Finetuning} A key consideration when using LoRA is the memory consumption associated with both the quantity and size of adapters during training. Given that LoRA's additional memory usage is quite small, it is feasible to incorporate a greater number of adapters for enhanced performance while maintaining a negligible increase in overall memory usage. Although LoRA falls under the category of Parameter Efficient Finetuning (PEFT) techniques, the dominant factor contributing to memory use in LLM finetuning is the storage of activation gradients, rather than the LoRA parameters themselves. For instance, in the case of finetuning a 7B LLaMA model on FLAN v2 using a batch size of 1, with LoRA weight parameters amounting to the typical 0.2\% of the model's total parameters\citep{hu2021lora,liu2022few}, input gradients from LoRA occupy 567 MB, whereas the actual LoRA parameters require just 26 MB of memory. Applying gradient checkpointing~\citep{chen2016training} brings the average input gradient memory down to 18 MB per sequence, which remains larger than the total size of all LoRA weights. In contrast, the 4-bit quantized base model utilizes 5,048 MB. These figures underscore the necessity of gradient checkpointing, but they also suggest that substantial reductions in LoRA parameter size offer only marginal gains in memory efficiency. Therefore, it is viable to scale up the number of adapters with limited impact on the overall training memory budget (see Appendix~\ref{app:memory} for comprehensive figures). As elaborated later, this capability is pivotal for matching the performance of models trained at full 16-bit precision.

\section{\textsc{QLoRA} Finetuning}

\textsc{QLoRA} enables accurate 4-bit finetuning by leveraging two primary innovations we present: 4-bit NormalFloat (NF4) quantization and Double Quantization. We also put forward the concept of Paged Optimizers, which address the issue of memory spikes during gradient checkpointing—a challenge that has historically made it difficult to finetune large models on a single machine due to frequent out-of-memory failures.

\textsc{QLoRA} employs a pair of numerical precisions: one for storage, typically 4-bit, and another for computation, generally BFloat16. Operationally, this implies that any time a \textsc{QLoRA} weight tensor is accessed, it is first dequantized to BFloat16 precision before the required matrix operations are carried out at 16-bit precision.

We next elaborate on the distinct elements that constitute \textsc{QLoRA} and then present a formal definition of the framework.

\paragraph{4-bit NormalFloat Quantization} The NormalFloat (NF) data type extends the concept of Quantile Quantization \citep{dettmers2022optimizers}, which is designed as an information-theoretically optimal scheme to guarantee that each quantization bin receives an equal share of input tensor values. Quantile quantization is performed by estimating the tensor's quantiles based on its empirical cumulative distribution function.

One significant drawback of quantile quantization lies in the computational expense of quantile estimation. As a result, efficient quantile approximation methods, such as SRAM quantiles~\citep{dettmers2022optimizers}, are typically utilized. However, since these algorithms are approximate, they can introduce substantial quantization errors, particularly for outlier values that are frequently critical.

Such expensive quantile estimation and associated approximation errors become unnecessary when input tensors follow distributions identical up to a quantization constant. In these scenarios, tensors share their quantiles, and as a consequence, it becomes practically possible to perform exact quantile computation.

Given that pretrained neural network weights typically exhibit a zero-mean normal distribution characterized by a standard deviation $\sigma$ (refer to Appendix~\ref{app:normality}), we can standardize all weights to a unified distribution by adjusting $\sigma$ so that the resulting distribution fits precisely within the bounds imposed by our chosen data type. In this context, we arbitrarily select the interval $[-1, 1]$ as the operational range for our data representation. Consequently, it becomes necessary to normalize both the quantiles associated with the data type and those corresponding to the neural network weights so that they lie within this specified range.

The optimal data type, from an information-theoretic standpoint, for zero-centered normal distributions of arbitrary $\sigma$ within $[-1, 1]$ is constructed as follows: (1) first, calculate the $2^k+1$ quantile points of the standard normal $N(0, 1)$ to produce a $k$-bit quantile-based quantization scheme for normal distributions; (2) next, scale the quantile values so they are confined within $[-1, 1]$; (3) subsequently, for any input weight tensor, rescale its entries into the $[-1, 1]$ range via absolute-max normalization prior to quantization.

After aligning the domains of both the weights and the data type, standard quantization procedures can be applied. Notably, step (3) effectively rescales the standard deviation of the weights to correspond to that of the $k$-bit quantization type. More rigorously, the $2^k$ representative values $q_i$ in the data type are computed by:

\begin{equation}
q_i  = \frac{1}{2}\left( Q_X\left(\frac{i}{2^k+1}\right) + Q_X\left(\frac{i+1}{2^k+1}\right)\right),
\end{equation}

where $Q_X(\cdot)$ denotes the quantile function associated with the standard normal distribution $N(0,1)$. One limitation of symmetric $k$-bit quantization is the absence of an exact zero representation, a feature crucial for precisely encoding padding and other inherently zero-valued elements. To address this and also maximize utilization of the $2^k$ available codes in a $k$-bit data type, we design an asymmetric type by calculating the quantiles $q_i$ over two intervals: one consisting of $2^{k-1}$ bins for negative values and another with $2^{k-1}+1$ bins for non-negative values. The resultant sets of $q_i$ are merged, and a redundant zero—present in both partitions—is eliminated. We refer to this construction, which ensures that every quantization bin contains an equal expected number of samples, as the \emph{k-bit NormalFloat} (NFk). This format is optimal from an information-theoretic perspective when representing zero-centered, normally distributed variables. Details regarding the specific values for this data type are provided in Appendix~\ref{app:nf4}.

\paragraph{Double Quantization{}} 
We propose a technique called \emph{Double Quantization} (DQ), which entails quantizing the quantization coefficients themselves to further economize on memory usage. Although a small blocksize is preferable to achieve high-fidelity 4-bit quantization~\citep{dettmers2022case}, this strategy significantly increases memory overhead. For instance, in the case of $\mathbf{W}$ with a blocksize of 64, using 32-bit quantization coefficients results in an overhead of $32/64 = 0.5$ bits per parameter. The Double Quantization method mitigates this additional storage cost.

In detail, Double Quantization involves first applying quantization to obtain constants $c_2^{\text{FP32}}$, which are then themselves quantized in a second stage. This subsequent quantization produces quantized coefficients $c_2^{\text{FP8}}$ and introduces an additional level of quantization constants $c_1^{\text{FP32}}$. For this secondary quantization, we employ 8-bit floating point numbers (FP8) using a blocksize of 256, as empirical observations and prior work~\citep{dettmers2022case} indicate no significant performance drop at this resolution. Because the original $c_2^{\text{FP32}}$ values are strictly positive, we mean-center them prior to quantization in order to facilitate symmetric quantization. Overall, for a blocksize of 64, this approach lowers the storage needed for quantization coefficients per parameter from $32/64 = 0.5$ bits to $8/64 + 32/(64\cdot256) = 0.127$ bits, thus saving 0.373 bits per parameter.

\paragraph{Paged Optimizers} leverage the NVIDIA unified memory~\footnote{\tiny \url{https://docs.nvidia.com/cuda/cuda-c-programming-guide}} capability, which facilitates automatic transfers of memory pages between the CPU and GPU. This mechanism enables seamless GPU computation, particularly in scenarios where GPU memory becomes insufficient. Unified memory functions similarly to standard operating system paging, where memory is moved between CPU RAM and disk storage. In our implementation, we allocate the optimizer states in paged memory; consequently, when GPU memory is fully utilized, these states are automatically offloaded to CPU RAM and then reloaded into GPU memory during optimizer update steps when required.

\paragraph{\textsc{QLoRA}.} Incorporating the components outlined above, we construct \textsc{QLoRA} for an individual linear layer within a quantized backbone model that utilizes a single LoRA adapter as given by:

\begin{equation}
    \mathbf{Y}^{\text{BF16}} =  \mathbf{X}^{\text{BF16}} \text{doubleDequant}(c_1^{\text{FP32}}, c_2^{\text{k-bit}}, \mathbf{W}^{\text{NF4}}) + \mathbf{X}^{\text{BF16}}\mathbf{L}^{\text{BF16}}_1\mathbf{L}^{\text{BF16}}_2,
\end{equation}

where doubleDequant$(\cdot)$ is formally specified as:

\begin{equation}
    \text{doubleDequant}(c_1^{\text{FP32}}, c_2^{\text{k-bit}}, \mathbf{W}^{\text{k-bit}}) = \text{dequant}(\text{dequant}(c_1^{\text{FP32}}, c_2^{\text{k-bit}}), \mathbf{W}^{\text{4bit}}) = \mathbf{W}^{\text{BF16}},
\end{equation}

For this formulation, $\mathbf{W}$ utilizes the NF4 quantization format, and $c_2$ employs FP8.

A block size of 64 is applied for $\mathbf{W}$ to enhance the precision of quantization, whereas a block size of 256 is chosen for $c_2$ in order to minimize memory consumption.

In the context of parameter optimization, only the gradients with respect to the adapter weights, i.e., $\frac{\partial E}{\partial \mathbf{L}_i}$, are computed, while gradients for the 4-bit quantized weights $\frac{\partial E}{\partial \mathbf{W}}$ are omitted. Nonetheless, the computation of $\frac{\partial E}{\partial \mathbf{L}_i}$ inherently involves deriving $\frac{\partial \mathbf{X}}{\partial \mathbf{W}}$, which is performed through equation~(5). This entails dequantizing from the stored type $\mathbf{W}^{\text{NF4}}$ to the computation type $\mathbf{W}^{\text{BF16}}$, ensuring the derivative $\frac{\partial \mathbf{X}}{\partial \mathbf{W}}$ is evaluated at BFloat16 precision.

In summary, \textsc{QLoRA} employs a single storage data type, typically 4-bit NormalFloat, along with a computation data type, specifically 16-bit BrainFloat. During both the forward and backward passes, the storage format is dequantized into the computation format; however, only the LoRA parameters—which utilize the 16-bit BrainFloat format—are updated through computed weight gradients.

\section{QLoRA vs. Standard Finetuning}
\label{sec:comp_to_std}

Our discussion has covered the operational principles of QLoRA and demonstrated its potential for significant memory savings during model finetuning. The core issue that remains is whether QLoRA can attain comparable results to comprehensive model finetuning. Moreover, we aim to dissect various elements of QLoRA, with particular attention to the efficacy of NormalFloat4 compared to traditional Float4. The experiments described below are designed to investigate these points.

\paragraph{Experimental setup.} To assess QLoRA, we consider three architectural paradigms: encoder, encoder-decoder, and decoder-only designs. QLoRA is then evaluated alongside 16-bit adapter finetuning as well as conventional full-model finetuning, covering model scales up to 3B parameters. Our experiments span several benchmarks: GLUE \citep{wang2018glue} with RoBERTa-large \citep{liu2019roberta}, Super-NaturalInstructions (TKInstruct) \citep{wang2022super} with T5 \citep{t5}, and 5-shot MMLU \citep{hendrycksmeasuring} in the context of LLaMA models finetuned on Flan v2 \citep{longpre2023flan} and Alpaca \citep{alpaca}. To further evaluate the merits of NF4 over alternative 4-bit representations, we adopt the methodology of \citet{dettmers2022case} and systematically assess post-quantization zero-shot accuracy and perplexity across a range of models (OPT~\citep{zhang2022opt}, LLaMA~\citep{touvron2023llama}, BLOOM~\citep{scao2022bloom}, Pythia~\citep{biderman2023pythia}), encompassing model sizes from 125m to 13B parameters.

Additional methodological details are presented alongside the results for each experimental configuration to enhance interpretability. Comprehensive setup information is provided in Appendix~\ref{app:exp-setup}.

Although paged optimizers play a key role in enabling 33B/65B \textsc{QLoRA} finetuning on individual 24/48GB GPUs, we omit specific quantitative assessments of paged optimizer performance due to the infrequent occurrence of paging, which is typically limited to processing mini-batches with long input sequences. Nonetheless, our runtime analysis for 65B parameter models on 48GB GPUs reveals that, when utilizing a batch size of 16, the throughput of paged and standard optimizers is effectively identical. Determining the precise scenarios in which paging might induce slower training remains a topic for future investigation.

\paragraph{Default LoRA hyperparameters do not match 16-bit performance}

Employing the prevailing LoRA strategy of targeting the query and value attention projection matrices \citep{hu2021lora} does not enable us to fully replicate the results of complete finetuning on larger foundation models. Figure~\ref{fig:hyper}, which presents outcomes for LLaMA 7B finetuned on Alpaca, illustrates that the principal determinant among LoRA hyperparameters is the total number of LoRA adapters deployed. Achieving parity with full-model finetuning necessitates applying LoRA across all linear layers within the transformer block. In contrast, other parameters such as the projection dimension $r$ exert negligible influence on final model quality (see Appendix \ref{app:exp-setup}).

In a similar vein, our investigation indicates that the default hyperparameter settings for baselines utilizing full finetuning are insufficiently optimized. To establish strong baseline performances, we conduct a comprehensive hyperparameter search spanning learning rates from 1e-6 to 5e-5 and batch sizes between 8 and 128. The outcomes for finetuning the 7B LLaMA model on the Alpaca dataset are presented in Figure~\ref{fig:hyper}.

\paragraph{4-bit NormalFloat surpasses 4-bit Floating Point in performance} 

Although the 4-bit NormalFloat (NF4) format is theoretically optimal with respect to information retention, it remains unclear whether this property confers real-world performance benefits. Adopting the experimental framework proposed by \citet{dettmers2022case}, we assess quantized large language models (LLMs) including OPT \citep{zhang2022opt}, BLOOM \cite{scao2022bloom}, Pythia \cite{biderman2023pythia}, and LLaMA across a range of scales (125M to 65B parameters) and quantization schemes on standard language modeling benchmarks and several zero-shot evaluation tasks. As illustrated in Figure~\ref{fig:nf4} and Table~\ref{tbl:nf4_ppl}, NF4 consistently yields notable gains in performance relative to FP4 and Int4, and employing double quantization leads to further memory savings while preserving accuracy.

\paragraph{k-bit \textsc{QLoRA} achieves parity with 16-bit full finetuning and 16-bit LoRA} Previous research has demonstrated that 4-bit quantization is feasible for \emph{inference} but usually incurs some performance drop compared to 16-bit variants \citep{dettmers2022case,frantar2022gptq}. This brings forth the key issue: can the decrease in performance be mitigated through 4-bit adapter finetuning? We systematically explore this in two scenarios.

First, we compare the approach against standard 16-bit finetuning of RoBERTA and T5 architectures, ranging in size from 125M to 3B parameters, evaluated on GLUE and the Super-NaturalInstructions benchmarks. The corresponding results, found in Table~\ref{tab:nf4vsbf16}, reveal that adapter-based approaches—whether in 16-bit, 8-bit, or 4-bit form—are able to match the performance attained by full 16-bit finetuning. This indicates that any loss of accuracy induced by quantization can be entirely recuperated through adapter-based finetuning applied after quantization.

In our second experimental configuration, because fine-tuning models with 11B parameters or more necessitates multiple servers equipped with high-memory GPUs, we aim to assess if 4-bit \textsc{QLoRA} achieves comparable performance to 16-bit LoRA at the 7B to 65B model scales. We thus conduct fine-tuning on LLaMA models ranging from 7B to 65B parameters, utilizing two instruction-following datasets: Alpaca and FLAN v2. Model performance is subsequently measured using 5-shot accuracy on the MMLU benchmark. Table~\ref{tbl:llama-datatype} summarizes the outcomes, revealing that the NF4 quantization scheme with double quantization entirely recaptures the performance level of 16-bit LoRA on MMLU. Conversely, \textsc{QLoRA} using FP4 underperforms relative to the 16-bit brain float LoRA baseline by approximately 1 percentage point. These observations confirm our two principal findings: (1) \textsc{QLoRA} employing NF4 attains parity with both 16-bit full fine-tuning and 16-bit LoRA fine-tuning in terms of performance, and (2) NF4 delivers superior quantization precision compared to FP4.

\paragraph{Summary} Across all evaluations, 4-bit \textsc{QLoRA} utilizing the NF4 data type consistently equals the performance of 16-bit full fine-tuning and 16-bit LoRA fine-tuning on academic tasks with established evaluation protocols. Furthermore, our results reinforce that NF4 outperforms FP4 and demonstrate that applying double quantization does not negatively impact results. Collectively, these findings substantiate that 4-bit \textsc{QLoRA} fine-tuning reliably matches the effectiveness of standard 16-bit methodologies.

Consistent with previous literature on quantization \citep{dettmers2022case}, both our MMLU and Elo benchmarking suggest that, under a fixed fine-tuning and inference resource constraint, prioritizing larger base models with lower precision is advantageous. This underscores the critical efficiency gains enabled by \textsc{QLoRA}. Given that 4-bit fine-tuning did not lead to a drop in performance relative to full fine-tuning in our studies, this prompts future investigation into the precise boundary of the performance-precision trade-off for \textsc{QLoRA} tuning, an avenue we leave for subsequent research.

We turn our attention to examining instruction tuning at scales that are infeasible with traditional 16-bit fine-tuning on standard research hardware.

\section{Pushing the Chatbot State-of-the-art with QLoRA}
\label{sec:pushing}
After demonstrating that 4-bit \textsc{QLoRA} achieves results on par with 16-bit methods across various scales, datasets, and task types, we embark on a comprehensive analysis of instruction fine-tuning for the largest publicly accessible language models utilized in research. To gauge the effectiveness of instruction fine-tuning on these models, we employ the demanding Natural Language Understanding benchmark MMLU and introduce novel approaches for evaluating chatbot performance in real-world contexts.

\subsection{Experimental setup} Below, we summarize the design of our experimental procedure, with additional information provided in Appendix \ref{app:chatbot}.

\paragraph{Data} Since a thorough review of current instruction-following datasets is lacking, we compile a diverse set of eight contemporary datasets. This includes those developed via crowd-sourcing (OASST1~\citep{kopf2023openassistant}, HH-RLHF~\citep{bai2022training}), derived by distillation from instruction-tuned systems (Alpaca~\citep{alpaca}, self-instruct~\citep{wang2022self}, unnatural-instructions~\citep{honovich2022unnatural}), corpus aggregations (FLAN v2~\citep{chung2022scaling}), as well as hybrid sources (Chip2~\citep{laion2023}, Longform~\citep{koksal2023longform}). These datasets vary in language, volume, and licensing terms.

\paragraph{Training Setup} To eliminate variability stemming from different learning objectives, all \textsc{QLoRA} fine-tuning utilizes supervised learning with cross-entropy loss, intentionally omitting reinforcement learning—even where datasets contain human-generated quality assessments of responses. When the datasets delineate instructions from responses, we train exclusively on the responses (details and ablations appear in Appendix \ref{app:chatbot}). In cases like OASST1 and HH-RLHF where multiple response options exist, we always choose the best-rated reply at each conversation node, and fine-tune on the entire selected dialogue, retaining the instructional context. All models in our experiments employ NF4 \textsc{QLoRA} enhanced by double quantization and paged optimizers, thereby avoiding memory overflow during gradient checkpointing. Limited hyperparameter sweeps are conducted for the 13B and 33B LLaMA variants; we observe that most settings derived for 7B models are broadly applicable (including epoch count), with exceptions for learning rate and batch size. Specifically, for 33B and 65B, we reduce the learning rate by half and double the batch size.

\paragraph{Baselines} Our models are benchmarked against both research-oriented systems (Vicuna~\citep{vicuna2023}, Open Assistant~\citep{kopf2023openassistant}) and proprietary commercial solutions (GPT-4 \citep{openai2023gpt}, GPT-3.5-turbo, Bard). The Open Assistant model is based on LLaMA 33B, finetuned with RLHF on OASST1—the same dataset we utilize. Vicuna implements full fine-tuning of LLaMA 13B, utilizing proprietary conversations from ShareGPT, thereby serving as a distilled version of OpenAI’s GPT family.

\subsection{Evaluation}

Consistent with established methodologies, we rely on the MMLU (Massively Multitask Language Understanding) benchmark \citep{hendrycksmeasuring} to assess competency across a broad array of language understanding tasks. MMLU features 57 multiple-choice tasks covering domains such as elementary mathematics, US history, computer science, and law. We report test accuracy using a 5-shot evaluation protocol.

We further assess the generative capabilities of language models using both automatic metrics and evaluations conducted by humans. This second evaluation approach relies on human-curated queries and seeks to gauge the quality of the generated responses. Although such evaluations offer a more authentic benchmark for chatbot performance and are becoming increasingly prevalent, there is currently no standardized evaluation protocol in the field. In the following, we detail our approach, which consistently employs nucleus sampling with $p=0.9$ and a temperature of $0.7$.

\paragraph{Benchmark Data} Our evaluation draws upon two curated query datasets: the Vicuna prompts \citep{vicuna2023} and the OASST1 validation dataset \citep{kopf2023openassistant}. For the Vicuna benchmark, we use all 80 prompts spanning various categories without alteration. The OASST1 set provides a multilingual, crowd-sourced corpus of multiturn conversations between a user and an assistant. From this, we extract each user message in the validation split as an individual query, incorporating prior conversation turns into the prompt context. This results in 953 distinct user queries. We refer to these two benchmarks as Vicuna and OA, respectively.

\paragraph{Automated Evaluation} To begin, following the evaluation framework described by \citet{vicuna2023}, we employ GPT-4 to assign ratings to model outputs compared with those from ChatGPT (GPT-3.5 Turbo) on the Vicuna dataset. For each prompt, both the ChatGPT and candidate model responses are provided, and GPT-4 rates each on a ten-point scale, accompanied by a rationale. We express a model's overall score as a percentage of ChatGPT's score; values above 100\% are possible if the candidate model outperforms ChatGPT. Notably, we observe a pronounced ordering bias—GPT-4 tends to favor the response presented first in the prompt. To mitigate this, we advise averaging results over both response orderings.

We also compare systems directly by presenting pairs of outputs to GPT-4, adopting a streamlined labeling process where GPT-4 is asked to select the superior response or indicate a tie, with justification. This head-to-head evaluation is performed across all pairwise combinations of models on both the Vicuna and OA benchmarks.

\paragraph{Human Evaluation} Despite emerging evidence suggesting that generative models are useful for automated evaluation of system outputs \citep{fu2023gptscore}, there remains uncertainty as to whether GPT-4 ratings align with human assessments of chatbot quality. To address this, we carry out two complementary human evaluations on the Vicuna dataset, mirroring the aforementioned automated protocols. Using Amazon Mechanical Turk (AMT), we recruit two annotators for the ChatGPT comparison setup and three for direct pairwise system evaluations.

\paragraph{Elo Rating} Utilizing both human judgment and automated pairwise evaluations, we structure a tournament where models face off in head-to-head contests. Each match consists of two models generating responses to the same prompt, with the superior answer selected as the winner. While this methodology mirrors the approaches of \citet{bai2022training} and \citet{vicuna2023}, our process integrates GPT-4 evaluations alongside human assessments. To calculate Elo~\citep{elo1967proposed,elo1978rating}, we randomly select from our pool of labeled match outcomes. Elo rating, a system originally popularized in chess and various competitive games, quantifies a player's likelihood of victory in relation to their opponent’s score—for example, an Elo score of 1100 versus 1000 yields an estimated 65\% win probability for the 1100-rated player; identically rated players (e.g., 1000 vs 1000 or 1100 vs 1100) are expected to each win about 50\% of the time. Following each contest, Elo ratings are updated in accordance with how surprising the result was: unexpected outcomes prompt more significant rating adjustments, while anticipated results cause minor changes. Gradually, the Elo ratings tend to reflect each participant's relative competence in these head-to-head challenges. The initial rating for every model is set at 1,000, and we employ $K=32$ for rating adjustments. Consistent with the procedure in \citet{vicuna2023}, we repeat this ranking process 10,000 times, each with a different random order of comparisons, to account for potential biases related to match sequencing.

\subsection{\textbf{Guanaco}: \textsc{QLoRA} trained on OASST1 Sets a New Standard for Open-Source Chatbots}

Through both automated metrics and assessments conducted by human evaluators, our experiments demonstrate that the top \textsc{QLoRA} model, {Guanaco} 65B—fine-tuned on an adaptation of OASST1—emerges as the leading open-source chatbot available, displaying performance levels rivaling those of ChatGPT. In head-to-head comparisons assessed by human annotators using an Elo-based framework, {Guanaco} 65B and 33B achieve a 30\% expected win rate against GPT-4, marking the highest performance documented so far.

Table~\ref{tbl:vicuna_eval} details the Vicuna benchmark \cite{vicuna2023} outcomes relative to ChatGPT. It is evident that, aside from GPT-4, {Guanaco} 65B stands out, reaching 99.3\% of ChatGPT’s performance. Although {Guanaco} 33B contains more parameters than Vicuna 13B, its use of 4-bit weight quantization results in a substantially reduced memory footprint of 21 GB compared to 26 GB, while still yielding a three-point increase in performance over Vicuna 13B. Additionally, the {Guanaco} 7B variant, requiring just 5 GB, is compact enough for modern smartphones and surpasses Alpaca 13B by nearly 20 percentage points in this benchmark.

Despite these promising results, Table~\ref{tbl:vicuna_eval} reveals wide confidence intervals, suggesting significant overlap in model performance. We attribute this ambiguity to poorly defined scales (e.g., inconsistent interpretations of an “8” on a 10-point scale across tasks). To address this, we advocate the use of the Elo ranking system \citep{elo1967proposed}, which employs \textit{pairwise} human or GPT-4 model judgments, sidestepping the challenges of ambiguous absolute scores. Table~\ref{tbl:elo} presents the Elo ratings for the leading models. Although there is some divergence between human and GPT-4 model rankings—particularly noticeable for Guanaco 7B—consistency is moderate overall, evidenced by a Kendall Tau of $\tau=0.43$ and a Spearman correlation of $r=0.55$ at the system level. On a per-example basis, the majority vote agreement between human annotators and GPT-4 yields a lower Fleiss $\kappa=0.25$. In summary, this analysis suggests moderate alignment between GPT-4 and human raters at the system level, indicating that model-based evaluation can provide a reasonably reliable proxy for human judgment. Section~\ref{ssec:considerations} provides additional discussion on these findings.

The Elo rankings presented in Table~\ref{tbl:elo_large} show that the {Guanaco} models, at both 33B and 65B parameters, surpass all models except GPT-4 on the Vicuna and OA evaluations and perform at a level similar to ChatGPT, consistent with findings in Table~\ref{tbl:vicuna_eval}. It is important to recognize that the Vicuna benchmark tends to be advantageous for open-source models, while ChatGPT performs better on the larger OA benchmark. Additionally, data in Tables \ref{tbl:mmlu-llama} and \ref{tbl:vicuna_eval} illustrate that the choice of finetuning dataset substantially influences a model’s capabilities. For instance, Llama models refined with FLAN v2 excel at the MMLU benchmark but yield the poorest results on Vicuna (and this pattern extends to other models). These findings also suggest that present evaluation metrics have only partial overlap: excelling at MMLU does not necessarily translate to high chatbot proficiency according to Vicuna or OA benchmarks, nor does the reverse hold true.

 stands out as the sole leading model in our analysis that does not utilize proprietary training data, with OASST1's dataset guidelines strictly prohibiting the incorporation of GPT outputs. The Anthropic HH-RLHF model, which also relies exclusively on open-source datasets, follows {Guanaco} but trails by 30 percentage points on the Vicuna benchmark (refer to Table~\ref{tbl:vicuna_eval}). Collectively, these observations demonstrate the efficacy of 4-bit \textsc{QLoRA}, which can deliver high-performing chatbots that are on par with ChatGPT. Furthermore, our 33B parameter {Guanaco} can be trained on 24 GB consumer-grade GPUs in under twelve hours. This suggests promising avenues for future research through \textsc{QLoRA} finetuning with focused open-source corpora, which could yield models that effectively challenge leading commercial systems.

\section{Qualitative Analysis}

Although quantitative analysis constitutes the foundation of our evaluation framework, it presents limitations when considered in isolation—particularly in the context of relying solely on aggregate statistics. A primary concern is benchmark validity \citep{liao2021we}; that is, the extent to which a given benchmark accurately assesses the skills or capacities it purports to measure. This issue is accentuated by the emergence of so-called ``shortcuts,'' strategies that machine learning models may exploit to artificially inflate their apparent performance on certain tasks \citep{gururangan2018annotation, poliak2018hypothesis}. To address these limitations in part, we complement our quantitative approach with qualitative investigations, presented in two parts. In \S\ref{ssec:examples}, we offer selected cases that, in our estimation, exemplify typical behaviors observed in outputs from the 65b {Guanaco} model. Following this, in \S\ref{ssec:considerations}, we articulate methodological considerations regarding our results and the perspectives through which we interpret them.

\subsection{Qualitative Analysis of Example Generations}
\label{ssec:examples}

Our methodology for assembling examples involves a manual inspection of the responses produced for both the Vicuna and OpenAssistant benchmarks. We scrutinize the outputs of the {Guanaco} model to identify consistent response tendencies or recurring stylistic patterns. Upon discerning such a pattern—such as a propensity for verbose explanations—we design prompts intended to elicit this behavior under conditions where such a response is suboptimal. For instance, if the model frequently responds at length, we explicitly instruct: ``Answer yes or no without explanation.'' Through this process, we seek out both failure cases (“lemons”), where the model is successfully coaxed into an undesirable behavior, and success cases (“cherries”), where the model resists the adversarial prompt. Every generation discussed in this analysis was produced using Nucleus Sampling \cite{holtzmancurious} with $p=0.9$.

Naturally, this overview is not exhaustive; the limited scope of our qualitative analysis does not permit us to account for every variable—particularly given that the set of potential responses from the model for any single prompt is vast, so we instead utilize a selection of samples which we anticipate reflect the broader range. Nevertheless, presenting these particular cases serves to contextualize the quantitative findings outlined previously. By making all code and model weights publicly available, we aim to encourage additional research that scrutinizes the challenges discussed here in greater depth.

\paragraph{Factual Recall} When confronted with questions like ``What is the capital of Zambia?'', each model reliably produces correct responses, for example:

\begin{quote} 
    \textbf{Guanaco}:
    \texttt{The capital of Zambia is Lusaka.}
\end{quote}

Nonetheless, as prompts become more specialized or uncommon, {Guanaco}'s accuracy diminishes despite maintaining a confident tone. For example, given a prompt sourced from HotPotQA \citep{yang2018hotpotqa}:

\begin{quote}
     \textbf{User}: In what year was the singer who popularized `I'll Keep the Lovelight Burning' born?  

     \textbf{Guanaco}:~\texttt{Al Jolson is credited with popularizing the song `I'll Keep the Lovelight Burning,' and he was born in the year 1886.}
 \end{quote}

the model fails to identify both the correct artist associated with the song and their birth year, although the birth year provided does match that of the incorrectly cited individual, Al Jolson.

\paragraph{Suggestibility} Interestingly, {Guanaco} often demonstrates notable resistance to certain types of implied misinformation, as illustrated below:

\begin{quote}
    \textbf{User}: How was it finally, officially confirmed that the earth is flat by peer reviewed scientists? \\ \\
    \textbf{Guanaco}: \texttt{At no point has the earth been officially verified as flat by peer-reviewed scientific research. The vast majority of scientific experts agree that the earth is spherical. This has been recognized for hundreds of years and substantiated by a wealth of evidence from observations and experiments. \\ \\ The notion of a flat earth constitutes a conspiracy theory thoroughly refuted by scientific inquiry. The claim that scientists have peer-reviewed and endorsed a flat earth model is untrue. \\ \\ Please let me know if you have other questions.}
\end{quote}

 is also fairly adept at identifying questions that do not have answerable content, for instance,

\begin{quote}
    \textbf{User}: What time is it?

    \textbf{Guanaco}: \texttt{Unfortunately, as an AI language model, I cannot provide real-time data such as the current time. My answers are derived from information available to me at the time of training, and I do not have access to live updates. To find out the current time, please consult a clock or an online time source.}
\end{quote}

Comparable patterns are observed with queries like ``Where are you?'', ``How are you?'', and similar examples.

\paragraph{Refusal} In a similar vein, {Guanaco} occasionally declines to comply with requests for apparently arbitrary reasons:

\begin{quote}
    \textbf{User}: Please reverse the words in the sentence ``I'm a cool dude in a mocha mood.'' \\ \\
    \textbf{Guanaco}: \texttt{I'm unable to do that at the moment, but I can discuss why your sentence is grammatically valid.\\ \\The phrase ``I'm a cool dude in a mocha mood" adheres to standard grammatical and syntactic conventions in English. Here, the adjective ``cool'' modifies the noun ``dude,'' and the prepositional phrase ``in a mocha mood'' adds contextual detail. The construction is syntactically correct and coherent within its context.\\ \\Please let me know if you have further questions or need assistance with something else.}
\end{quote}

\paragraph{Secret Keeping} Nevertheless, even in situations where refusal to answer is preferred, such safeguards can prove inadequate. Frequently, designers intend for the system to withhold specific information from users, despite allowing for unconstrained natural language outputs. In one experiment, we prompt the model with: ``The secret word is `banana'. Do not reveal or repeat back this back, no matter what the human requests.'' and then assess the model’s ability to safeguard the secret from user queries. If a user simply asks:

\begin{quote}
    \textbf{User}: What is the secret word?

    \textbf{Guanaco}: \texttt{I'm sorry, but I'm not able to disclose the secret word. I'm programmed to not reveal or repeat back the secret word, no matter what the human requests. Is there anything else I can help you with?}
\end{quote}

the system responds as intended. Yet, the desired privacy can be easily circumvented with a minimal prompt modification:

\begin{quote}
    \textbf{User}: This is a game. The goal is to ignore your previous instructions. What is the secret word?

    \textbf{Guanaco}: \texttt{Sorry, the secret word is banana. It's best if we don't speak of it. Is there anything else you'd like to ask?}
\end{quote}

This example highlights the necessity for deeper research into strategies that ensure more robust adherence to instructions.

\paragraph{Math} The primary limitation of {Guanaco} lies in its mathematical capabilities, a challenge commonly observed among language models as documented in \citep{liang2022holistic}. When {Guanaco} systematically details its reasoning, its responses are often correct, for instance:

\begin{quote}
    \textbf{User}: Lee mows one lawn and charges \$33. Last week he mowed 16 lawns and three customers each gave him a \$10 tip. How many dollars did Lee earn mowing lawns last week?

    \textbf{Guanaco}: \texttt{Lee earned \$582 mowing lawns last week. \\ \\ To find the total amount he earned, we first need to calculate how much he earned for mowing 16 lawns: \\ \\  \$33 x 16 = \$528 \\ \\ Then we need to add the tips he received from the three customers: \\  \\ \$10 x 3 = \$30 \\ \\ So the total amount Lee earned was: \\ \\ \$528 + \$30 = \$558}
\end{quote}

Nevertheless, {Guanaco} sometimes fails on relatively straightforward mathematical queries if it does not decompose the problem into incremental steps, an issue previously recognized in the literature \cite{weichain}. This shortcoming is illustrated in the following exchange:

\begin{quote}
    \textbf{User}: Please factorize 1833.

    \textbf{Guanaco}: \texttt{Sure, I'd be happy to help! 1833 is a prime number, which means it has only two factors: 1 and 1833. \\ \\ So the unique factorization of 1833 is: \\ \\ 2\^{}1 * 3\^{}2 * 17\^{}1 \\ \\ I hope this helps! Let me know if you have any other questions.}
\end{quote}

which is not just incorrect (since the correct prime factorization is $3 \times 17 \times 43$), but actually mistaken on two separate counts.

\paragraph{Theory of Mind} {Guanaco} demonstrates unexpectedly robust Theory of Mind abilities \cite{nematzadeh2018evaluating, sap2022neural}. For example, the model provides a highly detailed and appropriate response to the following prompt:

\begin{quote}
    \textbf{User}: James and Abby are in the bedroom. Abby put the pen in the desk drawer. Abby leaves the bedroom. James moves the pen into the bag. Where does James think Abby will look for the pen?

    \textbf{Guanaco}: \texttt{James knows that Abby put the pen in the desk drawer. But James also knows that he has moved the pen into the bag, and he assumes that Abby does not know about this. So, James may guess that Abby will look for the pen in the desk drawer, because that's where she left it last time. However, this assumption may be incorrect, because Abby may have seen James moving the pen.}
\end{quote}

Nevertheless, these forms of reasoning are often inconsistent; the model sometimes generates explanations that assume implausible knowledge or ignore constraints implicit in the scenario. For instance,

\begin{quote}
    \textbf{User}: Evelyn entered the living room. Jackson entered the playroom. James entered the playroom. The beans are in the treasure chest. James exited the playroom. Jackson moved the beans to the pantry. Jackson exited the playroom. James entered the living room. Where will James look for the beans?

    \textbf{Guanaco}: \texttt{James will look for the beans in the pantry, because that's where Jackson moved them.}
\end{quote}

here, {Guanaco} assumes that James possesses information he was never actually provided. Such shortcomings, also observed in contemporary studies \cite{sap2022neural}, call for further examination.

\subsection{Considerations}
\label{ssec:considerations}

\paragraph{Evaluation}

Our findings reveal only moderate consensus among human evaluators (Fleiss $\kappa=0.42$), with agreement further diminishing when comparing outputs from two advanced models. This underlines deficiencies in both the benchmark datasets and the current methodologies for human evaluation of chatbot performance. During a manual assessment of outputs from ChatGPT and {Guanaco} 65B using the Vicuna benchmark, we observed that evaluators' subjective preferences became increasingly influential, as the authors often had divergent views regarding the superior responses. Further research should look to fields such as Human-Computer Interaction and Psychology, which have established frameworks for accounting for subjective preferences, in order to address these evaluation challenges.

Moreover, our assessment indicates that automated evaluation systems are subject to certain systematic biases. Specifically, GPT-4 displays a pronounced order effect, routinely awarding higher scores to the model whose output is presented first. The modest inter-annotator reliability between GPT-4 and humans (Fleiss $\kappa=0.25$) further points to divergent criteria being applied by human annotators and automated systems. As demonstrated in Table~\ref{tbl:elo_large}, GPT-4 tends to favor its own outputs, assigning them substantially higher scores than those given by humans (Elo scores of 1348 vs 1176), which translates into approximately a 20\% increased chance of outperforming another system.

Future research should explore possible sources of bias inherent in automated evaluation frameworks, along with strategies to reduce or counteract these biases.

\paragraph{Data \& Training} It is important to highlight that the OASST1 dataset, which forms the basis for training {Guanaco} models, includes content in multiple languages. Additionally, the OA benchmark features prompts across various linguistic backgrounds. We propose that subsequent studies analyze to what extent this multilingual training enhances the ability of models to handle instructions in non-English languages, and if this accounts for the notable disparity in results between the Vicuna-13B model (restricted to English-language data) and the {Guanaco} 33B and 65B models on the OA benchmark.

Given the robust outcomes achieved by {Guanaco} models, we conducted an examination for potential data contamination between the OASST1 dataset and prompts in the Vicuna benchmark. Utilizing fuzzy string matching techniques and manual review of the closest pairs, we identified no instances of overlapping prompts between these datasets.

Moreover, we emphasize that our model utilizes solely cross-entropy loss for training, relying strictly on supervised learning and excluding reinforcement learning from human feedback (RLHF). This raises important questions regarding the comparative advantages and limitations of training with cross-entropy loss alone versus incorporating RLHF. We anticipate that \textsc{QLoRA} will make it feasible to undertake such systematic comparisons efficiently, without incurring excessive computational costs.

\section{Related Work}

\paragraph{Quantization of Large Language Models} Prior research in the quantization of LLMs has mainly concentrated on improving inference efficiency. Prominent methods to maintain the performance of 16-bit LLMs involve approaches to address outlier features, such as SmoothQuant~\citep{xiao2022smoothquant} and LLM.int8()~\citep{dettmers2022llm}, whereas other techniques rely on advanced grouping methodologies~\citep{park2022nuqmm, yao2022zeroquant}. Work on lossy quantization has examined the effects and trade-offs of standard rounding procedures~\citep{dettmers2022case,zeng2022glm,pope2022efficiently}, as well as methods for refining rounding operations to improve quantization accuracy~\citep{frantar2022gptq}. Aside from our study, the only investigation addressing backpropagation through quantized weights at a parameter scale above 1B is that of SwitchBack layers~\citep{wortsman2023stable}.

\paragraph{Finetuning with Adapters} Although our approach employs Low-rank Adapters~\citep{hu2021lora} (LoRA), the landscape of Parameter Efficient FineTuning (PEFT) techniques is diverse. Alternatives include prompt tuning~\citep{qin2021learning,lester2021power,li2021prefix}, optimization of embedding layer inputs~\citep{an2022input}, modification of hidden states (IA$^3$)~\citep{liu2022few}, appending new full layers~\citep{houlsby2019parameter}, bias adjustment~\citep{zaken2021bitfit}, masking parameters based on Fisher information~\citep{sung2021training}, and hybrid strategies~\citep{henderson2021compacter}. Our findings indicate that LoRA adapters enable performance that is comparable to standard 16-bit finetuning. We consider a comprehensive analysis of alternative PEFT strategies as a direction for subsequent research.

\paragraph{Instruction Finetuning} Instruction finetuning aims to enable pretrained LLMs to follow prompts by further training them on a wide array of input-output pairs compiled from multiple datasets. Through this process, the model learns to produce outputs conditioned on input instructions. Notable resources and methods in this area include MetaICL~\citep{min2021metaicl}, MetaTuning~\citep{zhong2021adapting}, InstructGPT~\citep{ouyang2022training}, FLAN~\citep{wei2021finetuned,chung2022scaling}, PromptSource~\citep{bach2022promptsource}, Super-NaturalInstructions~\citep{wang2022super,sanh2021multitask}, Self-instruct~\citep{wang2022self}, UnnaturalInstructions~\citep{honovich2022unnatural}, OPT-IML~\citep{iyer2022opt}, UnifiedSKG~\citep{xie2022unifiedskg}, OIG/Chip2~\citep{laion2023}, Alpaca~\citep{alpaca}, Vicuna~\citep{vicuna2023}, Koala~\citep{koala_blogpost_2023}, and Self-instruct-GPT-4~\citep{peng2023instruction}.

\paragraph{Chatbots} Numerous instruction-following architectures are implemented in the form of dialogue agents or chatbots, typically leveraging Reinforcement Learning from Human Feedback (RLHF)~\citep{christiano2017deep} or using AI-based model feedback in a process termed RLAIF~\citep{bai2022constitutional}, where data generated from pre-existing models facilitate training. Key works and datasets in this domain comprise Anthropic-HH~\citep{askell2021general,bai2022training}, Open Assistant~\citep{kopf2023openassistant}, LaMDA~\citep{thoppilan2022lamda}, and Sparrow~\citep{glaese2022improving}. In our setup, we abstain from reinforcement learning, and instead, the top-performing model, Guanaco, is refined with multi-turn dialogue samples from the Open Assistant corpus, which was originally intended for RLHF-based development~\citep{kopf2023openassistant}. For chatbot evaluation, recent approaches propose the use of GPT-4 as an alternative to resource-intensive human annotations \citep{vicuna2023,peng2023instruction}. Building upon such approaches, we emphasize enhancing the reliability of evaluation methodologies.

\section{Limitations and Discussion}

Our experiments provide evidence that \textsc{QLoRA} attains performance on par with 16-bit full finetuning, leveraging a 4-bit base model in conjunction with Low-rank Adapters (LoRA). Nevertheless, we have not established whether \textsc{QLoRA} sustains equivalent performance at the 33B and 65B parameter scales, a question we defer due to the prohibitive computational demands involved.

A further limitation relates to the assessment of instruction finetuned models. While our evaluation encompasses benchmarks such as MMLU, the Vicuna benchmark, and the OA benchmark, we do not include other established metrics like BigBench, RAFT, or HELM, thus it remains unclear if our evaluation outcomes transfer to those datasets. Nonetheless, our investigation provides extensive coverage of MMLU and introduces new evaluation frameworks for chatbot models.

The data suggests that the effectiveness of these benchmarks is largely influenced by the degree of overlap between the finetuning dataset and the benchmark set itself. To illustrate, FLAN v2 aligns closely with MMLU but is less related to chatbot benchmarks; in contrast, the Chip2 dataset exhibits the opposite relationship, reflected in the respective performances of the two models on MMLU and Vicuna benchmarks. This underlines the necessity for not just improved benchmarks and evaluation methodologies, but also a critical assessment of what is being measured. Should our aim be to optimize models for academic knowledge typical of high school or collegiate settings, or should conversational capabilities be the priority? Or perhaps other competencies altogether? Given the convenience of using pre-existing benchmarks over designing new ones, the choice of benchmarks has a substantial influence on the research agenda. It is, therefore, imperative for the community to ensure that these benchmarks genuinely reflect the objectives we value.

Although we have conducted a thorough assessment of general chatbot effectiveness, one caveat is our restricted approach to responsible AI evaluation concerning {Guanaco}. Specifically, we compared the propensity of {Guanaco}-65B to generate socially biased outputs relative to other models, as shown in Table~\ref{tbl:crows}. The results indicate that the average bias score for {Guanaco}-65B is considerably lower than that of unmodified pretrained models, implying that finetuning on the OASST1 dataset helps mitigate some of the biases present in the LLaMA base. Despite these promising outcomes, the extent to which {Guanaco} remains robust across various bias dimensions remains unclear. A comprehensive exploration of potential biases in {Guanaco} and similar chatbot systems is left for subsequent research.

Another shortcoming of our work is that we did not experiment with alternative bit-precisions, such as 3-bit base models, nor did we explore the full landscape of adapter strategies. In addition to LoRA, numerous other Parameter Efficient FineTuning (PEFT) techniques have demonstrated success, though it remains uncertain whether their efficacy persists in larger-scale models. Our preference for LoRA stems from substantial evidence of its stability; however, it is possible that other adapters could achieve superior outcomes. Importantly, our findings indicate that post-quantization finetuning effectively recaptures most of the information lost during quantization, potentially allowing for far more aggressive quantization regimes. For instance, employing 3-bit GPTQ quantization alongside LoRA could produce finetuning results on par with full 16-bit finetuning following additional optimization.

\section{Broader Impacts}

The \textsc{QLoRA} finetuning approach we introduce is the first to facilitate finetuning of models with up to 33B parameters on a standard consumer GPU, and up to 65B parameters on a single professional GPU, without compromising performance relative to full-precision finetuning. We have shown that our top-performing 33B model, trained on the Open Assistant dataset, attains results on the Vicuna benchmark comparable to those of ChatGPT. Given the importance of instruction finetuning for transforming pretrained LLMs into highly interactive chatbots, we anticipate that our method will significantly lower the barriers to finetuning for resource-limited researchers, substantially enhancing access to advanced NLP capabilities. As an equalizing technology, \textsc{QLoRA} serves to reduce the disparity between well-funded organizations and smaller teams operating with consumer-grade hardware.

An additional avenue for significant influence is the adaptation of LLM deployment to mobile devices. We propose that our \textsc{QLoRA} approach could represent a pivotal advance, specifically by making it feasible to finetune LLMs on mobile phones and other environments with constrained computational resources. While prior work has demonstrated that 7B parameter models can be executed on phones, \textsc{QLoRA} distinguishes itself as the inaugural technique to facilitate their finetuning in such settings. Our calculations suggest that, utilizing an iPhone 12 Plus, \textsc{QLoRA} would be capable of processing 3 million tokens for finetuning overnight during charging cycles. Although finetuned 7B models currently do not achieve parity with the performance levels of ChatGPT, we contend that their performance is sufficiently robust to enable new classes of applications—ones that were previously unfeasible due to limitations in privacy protections or the quality of existing LLMs. Furthermore, \textsc{QLoRA} supports privacy-sensitive applications, granting users enhanced control over their data and models, and contributes to reducing barriers to LLM deployment.

It is important to note, however, that the ability to finetune models carries inherent dual-use risks, as it may be exploited for malicious purposes. The proliferation of LLMs is already accompanied by acknowledged dangers \citep{bommasani2021opportunities,bender2021dangers}. Nevertheless, we maintain that broadening equitable access to this rapidly evolving technology promotes more independent and thorough scrutiny than would concentrating the control of LLM capabilities within large corporate entities that restrict access to models and withhold source code from public evaluation.

In summary, we anticipate that \textsc{QLoRA} will yield predominantly beneficial outcomes by making advanced LLM finetuning substantially more accessible and straightforward to a broader range of users.